% SI Text 3: extended results. Full treatment of four results the main text states in
% one sentence (Figs. S2 to S5). Condensed 2026-09-29: main-text quotations removed
% (the claim-to-section table at the top of the SI maps them), footnotes folded into
% text, restated arguments cut.
\section{Extended results}\label{si:extended}

\subsection*{The value of targeting decreases in the budget}
The value of targeting is the expected research output the optimal allocation adds over equal division of the same budget. By Corollary~\ref{cor:targeting-vanishes} of \ref{app:gap-rule} it vanishes as the budget grows, for every production function under which research output has a ceiling set by capability; that it also decreases throughout is a numerical finding, which held in every population we sampled (\ref{app:gap-rule}). Fig.~\ref{fig:targeting-value} shows it for one population, as a fraction of uniform funding's gain over no funding. The same dependence on the budget sets how much research output seed grants and lotteries give up and when the informed funder concentrates rather than spreads its grants (main text).

% Figure 2: the value of targeting decreases in the budget.
% Computed exactly from Proposition 1 (no simulation): population of n = 75 with
% K, R0 ~ Pareto(alpha = 2, min 1), rho = 0, A = 1, seed 659 (numpy default_rng);
% value of targeting = (f_opt - f_unif) / (f_unif - f_0), in percent;
% b = budget / total baseline resources. Data: fig2-curve.dat, generated by
% verify_gap_rule.py's companion snippet (see state.md, 2026-09-05).
\begin{figure}[tp]
\centering
\begin{tikzpicture}
\begin{axis}[
  width=0.82\textwidth, height=5.4cm,
  axis lines*=left,
  xmin=0, xmax=2, ymin=0, ymax=175,
  xtick={0, 0.5, 1, 1.5, 2},
  ytick={0, 50, 100, 150},
  xlabel={budget scale $b$},
  ylabel={value of targeting (\%)},
  ylabel style={font=\small}, xlabel style={font=\small},
  tick label style={font=\small},
  axis line style={line width=0.4pt},
  tick style={line width=0.4pt, black},
  clip=false,
]
\addplot[black, line width=0.9pt, smooth] table {fig2-curve.dat};
\end{axis}
\end{tikzpicture}
\caption{The value of targeting against the budget. The optimal allocation's gain over uniform funding, as a fraction of uniform funding's gain over no funding, as the budget scale $b$ grows, for one population drawn from the default distributions. Corollary~\ref{cor:targeting-vanishes} proves that the value vanishes as the budget grows; that it falls throughout is a numerical finding (\ref{app:gap-rule}).}
\label{fig:targeting-value}
\end{figure}


For a funder allocating with records and review rather than with complete information, the value of targeting relative to uniform funding's gain over no funding decreases from 1.17 to 0.72 across budget scales $b = 0.2$ to $2$ with heavy-tailed capability ($\alpha_K = 1.3$) and reliable review ($\tau_K = 0.3$), and from 0.74 to 0.25 in the intermediate field with the default signal ($\alpha_K = 2$, $\tau_K = 1$); two rounds, 200 simulated populations per point, the populations of Fig.~4\textit{A}. Giving up targeting loses little research output under the same conditions that make review worth little: where capability is evenly spread, the budget is large, or review is uninformative.

\subsection*{Overtrusting review}
The results on review's value require the funder to weight review at its true informativeness. A funder that overtrusts review, treating a noisy signal as reliable, can do worse than one that ignores review altogether: with moderately spread capability, overtrust turns review's value negative; with heavy-tailed capability it reduces the value without reversing it (Fig.~\ref{fig:overtrust}). A funder that overtrusts tenfold (true $\tau_K = 3$, belief $0.3$) loses more than review's full value at the intermediate capability distribution, whereas a funder that undertrusts tenfold (true $\tau_K = 0.3$, belief $3$) retains over half of it. Overtrusting review loses more research output than undertrusting it.

% Figure 6 (S6): the value of review under misweighted trust.
% y = review's value (S5 - S4) as percent of no-funding output; x = the funder's
% believed noise (log scale); true noise tau_K = 3 (dotted rule). Two capability
% distributions. Error bars: one standard error, 200 seeds per cell.
% Data: fig6-overtrust.dat, from sweep_results/D_misspecified_trust (hash 21b0d9a).
\begin{figure}[tp]
\centering
\begin{tikzpicture}
\begin{axis}[
  width=0.82\textwidth, height=6.0cm,
  axis lines*=left,
  xmode=log,
  xmin=0.1, xmax=10, ymin=-5, ymax=45,
  xtick={0.1, 0.3, 1, 3, 10},
  xticklabels={0.1, 0.3, 1, 3, 10},
  ytick={0, 10, 20, 30, 40},
  xlabel={believed noise $\tau_K$},
  ylabel={value of review (\% of the gain from funding)},
  ylabel style={font=\small}, xlabel style={font=\small},
  tick label style={font=\small},
  axis line style={line width=0.4pt},
  tick style={line width=0.4pt, black},
  clip=false,
]
\addplot[gray, dashed, line width=0.4pt, domain=0.1:10] {0};
\draw[gray, dotted, line width=0.5pt] (axis cs:3,-2) -- (axis cs:3,20);
\node[font=\small, gray, anchor=west] at (axis cs:3.15,18) {true noise};
\addplot[black, line width=0.7pt, error bars/.cd, y dir=both, y explicit]
  table[x=belief, y=p13, y error=s13] {fig6-overtrust.dat};
\addplot[black, line width=0.7pt, error bars/.cd, y dir=both, y explicit]
  table[x=belief, y=p20, y error=s20] {fig6-overtrust.dat};
\node[font=\small, anchor=west] at (axis cs:11,25.5) {$\alpha_K = 1.3$};
\node[font=\small, anchor=west] at (axis cs:11,3.9)  {$\alpha_K = 2$};
\end{axis}
\end{tikzpicture}
\caption{Overtrusting review loses more research output than undertrusting it. The value of review as a percent of the research output that a correctly calibrated review-informed funder adds over no funding, as the funder's believed signal noise varies while the true noise is $\tau_K = 3$ (dotted rule); bars are one standard error. At the default capability distribution, overtrust takes the value below zero. Settings: \ref{app:simulation}.}
\label{fig:overtrust}
\end{figure}


\subsection*{Grant size, learning, and the schedule}
How much a funder can learn from its own grants depends on their size. A researcher on a small grant produces research output nearly proportional to the grant and nearly independent of capability, so a funder whose grants are small learns almost nothing it can exploit later: in a field whose researchers have few resources, with no review signal, such a funder gains almost nothing over uniform funding. As grants approach the scale of capability itself, research output separates the capable from the rest, and the funder recovers most of the value of discrimination, with timing that spends a small share early, so that research output can be observed, and the rest once the funder has learned who is capable (Fig.~\ref{fig:depth-schedule}). In a community whose researchers have few resources, with heavy-tailed capability and no review signal, over six rounds, the Bayesian funder's gain over uniform funding increases from a third of one percent of uniform funding's research output at the default grant size to nine percent at six times that size and ten percent at twelve times, roughly the gain a reliable review signal delivers at the default size.

The shift of spending toward later rounds comes from the research output researchers produce whether or not they are funded: that research output compounds and reveals capability while the funder waits. Growth from grant-added research output alone would shift spending slightly toward earlier rounds, since earlier grants have longer to compound. With growth from unfunded research output switched off, the center of mass of spending is 0.48 to 0.49, below even; with growth from grant-added research output switched off, it is 0.53 to 0.68 (five rounds).

% Figure (S7): the optimal schedule at standard and deep funding.
% y = share of the budget spent in each round by the forward funder (S8); T = 6;
% poverty corner (r_min = 0.001), heavy tail (alpha_K = 1.3), no review signal
% (tau_K = 100), negligible compounding (eps = 1e-4), budget scaled against
% capability (budget_ref = "K"). b = 0.5: 50 seeds; b = 3: 24 seeds (max SE 0.006);
% round-1 share at b = 3 matches the canonical 200-seed run (0.104 vs 0.110).
% Data: fig7-schedule.dat, generated by for-claude/verify_s7.R (schedule addendum).
\begin{figure}[tp]
\centering
\begin{tikzpicture}
\begin{axis}[
  width=0.66\textwidth, height=5.4cm,
  axis lines*=left,
  xmin=1, xmax=6, ymin=0.08, ymax=0.20,
  xtick={1, 2, 3, 4, 5, 6},
  ytick={0.10, 0.15, 0.167, 0.20},
  yticklabels={0.10, 0.15, even, 0.20},
  xlabel={round},
  ylabel={share of budget},
  ylabel style={font=\small}, xlabel style={font=\small},
  tick label style={font=\small},
  axis line style={line width=0.4pt},
  tick style={line width=0.4pt, black},
  clip=false,
]
\addplot[black, line width=0.7pt, mark=*, mark size=1.2pt] table[x=round, y=b05] {fig7-schedule.dat};
\addplot[black, line width=0.7pt, mark=o, mark size=1.4pt] table[x=round, y=b3] {fig7-schedule.dat};
\node[font=\small, anchor=west] at (axis cs:6.15,0.167) {default grant size ($b = 1$)};
\node[font=\small, anchor=west] at (axis cs:6.15,0.183) {large grants ($b = 6$)};
\end{axis}
\end{tikzpicture}
\caption{With large grants the funder spends a small share first and the rest once it has observed research output. The share of the budget spent in each of six rounds, in a field whose researchers have few resources, with heavy-tailed capability and no review signal, at the default grant size (filled) and at six times that size (open). Settings: \ref{app:simulation}.}
\label{fig:depth-schedule}
\end{figure}


\subsection*{Concentration and the production function}
One might hope that the production function settles whether funding should be concentrated among a few researchers or spread across many: the harder it is to substitute resources for capability, the more the money should go to the few researchers who can use it. Across the family of production functions of \ref{app:technology}, from Cobb-Douglas to Leontief, it does not (Fig.~\ref{fig:concentration}). With a tight budget in a heavy-tailed field, stronger complementarity does concentrate the review-informed funder's grants: the Gini coefficient of its grants rises from 0.92 to 0.95. With an ample budget, stronger complementarity spreads them, from 0.35 to 0.21 in the intermediate field and from 0.63 to 0.45 in a heavy-tailed field, a pattern consistent with a funder that, once the budget allows, brings every researcher's resources up toward the level their capability can use. Where capability is evenly spread and the budget ample, the relation is not even monotone: the Gini coefficient is 0.16 at Cobb-Douglas, peaks at 0.196 near $\gamma = -6$, and falls to 0.185 at $\gamma = -12$, 3.5 standard errors above the $\gamma = -12$ level. How widely funds should be spread is set jointly by the production function, the budget, and the field's capability distribution; no one of the three decides it alone.

% Figure (S8): how concentrated the informed allocation is, across the family.
% y = Gini coefficient of the review-informed funder's grants (single round);
% x = CES exponent gamma (0 = Cobb-Douglas, -1 = our harmonic form, limit =
% Leontief, plotted as detached marks at the left edge). Three field-budget pairs.
% Ground: sigma_tierB_concentration (+ leontief endpoint, + 400-seed refinement for
% the evenly spread curve), 200 seeds per cell, re-derived in-session (verify_s9.R).
% Leontief cells carry the kink caveat (ties broken by fill order, n_steps = 800).
\begin{figure}[tp]
\centering
\begin{tikzpicture}
\begin{axis}[
  width=0.82\textwidth, height=6.0cm,
  axis lines*=left,
  xmin=-14.8, xmax=0, ymin=-0.16, ymax=0.06,
  xtick={-14, -12, -6, -3, -1, 0},
  xticklabels={Leontief, $-12$, $-6$, $-3$, $-1$, $0$},
  ytick={-0.15, -0.10, -0.05, 0, 0.05},
  yticklabels={$-0.15$, $-0.10$, $-0.05$, 0, 0.05},
  xlabel={CES exponent $\gamma$ (complements $\leftarrow$, substitutable $\rightarrow$)},
  ylabel={change in grant concentration (Gini)},
  ylabel style={font=\small}, xlabel style={font=\small},
  tick label style={font=\small},
  axis line style={line width=0.4pt},
  tick style={line width=0.4pt, black},
  clip=false,
]
\draw[gray, dotted, line width=0.5pt] (axis cs:-1,-0.16) -- (axis cs:-1,0.06);
\node[font=\scriptsize, gray, anchor=south] at (axis cs:-1,0.061) {our form};
\addplot[black, line width=0.7pt, mark=*, mark size=1.1pt] table[x=g, y=dgini] {fig13-heavy01.dat};
\addplot[black, line width=0.7pt, mark=o, mark size=1.3pt] table[x=g, y=dgini] {fig13-def1.dat};
\addplot[black, line width=0.7pt, mark=triangle*, mark size=1.4pt] table[x=g, y=dgini] {fig13-even1.dat};
\addplot[black, only marks, mark=*, mark size=1.1pt] coordinates {(-14,0.031)};
\addplot[black, only marks, mark=o, mark size=1.3pt] coordinates {(-14,-0.147)};
\addplot[black, only marks, mark=triangle*, mark size=1.4pt] coordinates {(-14,0.009)};
\node[font=\scriptsize, anchor=north] at (axis cs:-8.5,0.014) {heavy-tailed, tight};
\node[font=\scriptsize, anchor=south] at (axis cs:-10.5,-0.047) {default, ample};
\node[font=\scriptsize, anchor=south] at (axis cs:-7.5,0.042) {evenly spread, ample};
\end{axis}
\end{tikzpicture}
\caption{Complementarity concentrates funding only when the budget is small. The change in the Gini coefficient of the review-informed funder's grants, relative to its Cobb-Douglas level, as the production function moves from Cobb-Douglas ($\gamma = 0$) toward Leontief (detached marks at the left edge), for three field-and-budget combinations. Settings: \ref{app:simulation}.}
\label{fig:concentration}
\end{figure}

